% ---------------------------------------------------------------------------
% main.tex
% Predicting the 2026 FIFA World Cup champion from the quarterfinal stage:
% a conditional Poisson score model with Monte Carlo simulation.
%
% Compiles with pdfLaTeX on Overleaf. No external images, no shell-escape.
% Companion files: algorithm.tex, model_notes.tex, references.bib,
%                  data_template.csv, simulate_worldcup.py, README.md
% ---------------------------------------------------------------------------
\documentclass[11pt]{article}

\usepackage[margin=2.5cm]{geometry}
\usepackage{amsmath}
\usepackage{amssymb}
\usepackage{booktabs}
\usepackage{graphicx}   % listed for completeness; no external images are used
\usepackage{algorithm}
\usepackage{algpseudocode}
\usepackage{hyperref}
\hypersetup{colorlinks=true, linkcolor=blue, citecolor=blue, urlcolor=blue}

\newcommand{\Prob}{\mathbb{P}}
\newcommand{\Expect}{\mathbb{E}}
\newcommand{\Ind}{\mathbf{1}}

\title{Predicting the 2026 FIFA World Cup Champion from the Quarterfinal Stage:\\
A Conditional Poisson Score Model with Monte Carlo Simulation}
\author{Tianyue Shao\\
\normalsize National University of Singapore}
\date{July 2026}

\begin{document}

\maketitle

\begin{abstract}
\noindent
We develop a fully specified probabilistic model for predicting the champion of
the 2026 FIFA World Cup \emph{conditional} on the tournament having reached the
quarterfinal stage. Goals in each remaining match follow independent Poisson
distributions whose log-rates combine team attack and defence parameters with
Elo-rating, in-tournament form, and rest-day covariates. Because every
remaining match is a knockout tie, draws after regular time are resolved by a
logistic tie-break driven by the Elo difference, which yields a closed-form
advancement probability for every ordered pair of teams. Championship
probabilities for the eight quarterfinalists are then estimated by Monte Carlo
simulation of the remaining seven matches. We describe how the match-level
probabilities are computed, provide pseudocode and a reference Python
implementation, and discuss interpretation, limitations, and extensions,
including expected goals, injuries, betting odds, and Bayesian parameter
uncertainty. All model inputs are estimated from real data through 7~July
2026: Elo ratings are replicated from the full international match record
(1872--2026), attack and defence parameters and covariate coefficients are
fitted by time-decayed Poisson regression on 2016--2026 internationals, and
the tie-break slope is fitted to 681 historical penalty shoot-outs. The
resulting probabilities are genuine model-based estimates, conditional on the
model's simplifying assumptions; data provenance, processing rules, and exact
reproduction steps are documented in the appendix.
\end{abstract}

\section{Introduction}
\label{sec:intro}

With the quarterfinal line-up of the 2026 FIFA World Cup decided, exactly eight
teams can still win the title and only seven matches remain: four
quarterfinals, two semifinals, and the final. This note formalizes the
question every observer now asks---\emph{given} that a team has reached the
quarterfinals, what is its probability of winning the World Cup?---and answers
it with a transparent, reproducible pipeline consisting of (i) a Poisson model
for match scores, (ii) covariate-driven scoring rates, (iii) a logistic
tie-break for drawn knockout matches, and (iv) Monte Carlo simulation of the
remaining bracket.

Independent Poisson models have been the workhorse of association-football
score prediction since Maher \cite{maher1982modelling}. Dixon and Coles
\cite{dixon1997modelling} refined the approach with an adjustment for the
dependence of low-scoring outcomes and with time-decayed estimation, and
bivariate extensions are analysed by Karlis and Ntzoufras
\cite{karlis2003analysis}. Aggregate team strength is convenient
to summarize through Elo-style ratings \cite{elo1978rating}, and
simulation-based forecasting of knockout tournaments is standard practice in
the sports-analytics literature \cite{koning2003simulation,groll2019hybrid}.
The model
presented here deliberately stays within this classical framework so that
every component is easy to estimate, audit, and extend.

Two design choices deserve emphasis. First, the prediction is
\emph{conditional}: the group stage and all earlier knockout rounds have been
played, so teams eliminated before the quarterfinals are excluded from the
sample space, and the realized bracket---who plays whom, and which semifinal
awaits each winner---is treated as known
(Section~\ref{sec:conditional}). Second, although only seven matches remain,
we estimate championship probabilities by Monte Carlo simulation rather than
by exact enumeration; the reasons, chiefly extensibility and richer outputs,
are discussed in Section~\ref{sec:mc} and Appendix~\ref{app:whymc}.

The note is organized as follows. Section~\ref{sec:conditional} defines the
conditional prediction problem and the bracket.
Section~\ref{sec:poisson} specifies the Poisson score model.
Section~\ref{sec:matchprob} derives match-level win and draw probabilities,
and Section~\ref{sec:draws} treats extra time and penalties.
Section~\ref{sec:mc} presents the Monte Carlo scheme, whose pseudocode appears
in Section~\ref{sec:algorithm}. Sections~\ref{sec:data}
and~\ref{sec:output} document the data inputs and the interpretation of the
output, and Section~\ref{sec:limitations} discusses limitations.
Appendix~\ref{app:notes} collects supplementary modelling notes. A reference
implementation is provided in the companion files \texttt{fit\_parameters.py}
(parameter estimation), \texttt{simulate\_worldcup.py} (tournament
simulation), and \texttt{data\_template.csv} (estimated team covariates).

\section{The Conditional Prediction Problem}
\label{sec:conditional}

Let
\[
\mathcal{T} = \{\text{France, Morocco, Spain, Belgium, Norway, England,
Argentina, Switzerland}\}
\]
denote the set of quarterfinalists, and let $C \in \mathcal{T}$ denote the
eventual champion. Write $\mathcal{Q}$ for the conditioning information: the
identity of the eight quarterfinalists together with the realized bracket,

\begin{itemize}
  \item \textbf{QF1:} France vs.\ Morocco,
  \item \textbf{QF2:} Spain vs.\ Belgium,
  \item \textbf{QF3:} Norway vs.\ England,
  \item \textbf{QF4:} Argentina vs.\ Switzerland,
  \item \textbf{SF1:} winner(QF1) vs.\ winner(QF2),
  \item \textbf{SF2:} winner(QF3) vs.\ winner(QF4),
  \item \textbf{Final:} winner(SF1) vs.\ winner(SF2).
\end{itemize}

The target of inference is the vector of conditional championship
probabilities
\begin{equation}
p_t \;=\; \Prob\bigl(C = t \,\big|\, \mathcal{Q}\bigr),
\qquad t \in \mathcal{T}.
\label{eq:target}
\end{equation}

Three properties follow directly from the conditioning. First, any team
eliminated before the quarterfinals is excluded from the sample space, so its
championship probability is zero \emph{by construction}; the model never
simulates such teams. Second, the probabilities in \eqref{eq:target} sum to
one over the eight quarterfinalists, because exactly one of them will lift the
trophy. Third, everything that happened before the quarterfinals---group-stage
results, earlier knockout rounds, accumulated minutes---enters the model only
through the team covariates measured at the quarterfinal snapshot (Elo rating,
tournament form, rest days; Section~\ref{sec:data}). Consequently, the numbers
produced here are \emph{not} comparable to pre-tournament title odds, which
average over the many brackets and line-ups that never materialized.

\section{Poisson Score Model}
\label{sec:poisson}

Consider a remaining match between teams $i$ and $j$, and let $G_i$ and $G_j$
denote the goals they score in regular time (90 minutes). We assume
conditionally independent Poisson counts,
\begin{equation}
G_i \sim \mathrm{Poisson}(\lambda_{ij}),
\qquad
G_j \sim \mathrm{Poisson}(\lambda_{ji}),
\label{eq:poisson}
\end{equation}
with a log-linear model for the scoring rates:
\begin{equation}
\log \lambda_{ij}
\;=\;
\alpha \;+\; a_i \;-\; d_j
\;+\; \gamma \,(E_i - E_j)
\;+\; \delta \,(F_i - F_j)
\;+\; \rho \,(R_i - R_j),
\label{eq:lograte}
\end{equation}
and symmetrically for $\lambda_{ji}$ with the roles of $i$ and $j$ exchanged.
The ingredients of \eqref{eq:lograte} are:
\begin{itemize}
  \item $\alpha$: a baseline log scoring rate, common to all teams, which
        absorbs the average goal level of knockout football at a neutral
        venue (there is no home-advantage term because World Cup knockout
        matches are treated as neutral-site fixtures);
  \item $a_i$: the \emph{attack strength} of team $i$ (larger values mean
        more goals scored, all else equal);
  \item $d_j$: the \emph{defence parameter} of team $j$. With the sign
        convention of \eqref{eq:lograte}, a larger $d_j$ \emph{reduces} the
        opponent's scoring rate, so $d_j$ measures defensive
        \emph{strength}. (Some authors instead add a defensive-weakness
        parameter $d_j^{\mathrm{w}}$ with a plus sign; the two conventions
        are equivalent via $d_j^{\mathrm{w}} = -d_j$.);
  \item $E_i$: the Elo rating of team $i$ entering the quarterfinal stage;
  \item $F_i$: a measure of \emph{tournament form}, e.g.\ a composite score
        of performances in this World Cup so far;
  \item $R_i$: rest days before the match (a fatigue proxy);
  \item $\gamma, \delta, \rho$: coefficients converting the covariate
        differences into log-rate effects.
\end{itemize}

Because only rating, form, and rest \emph{differences} enter
\eqref{eq:lograte}, the model is invariant to common shifts in the covariates.
The attack and defence parameters are identified only up to additive
constants (a shift of all $a_i$ can be absorbed into $\alpha$, and similarly
for $d_i$), so in estimation one imposes the usual sum-to-zero constraints
$\sum_i a_i = \sum_i d_i = 0$.

In practice the parameters
$\theta = (\alpha, \{a_i\}, \{d_i\}, \gamma, \delta, \rho)$ are estimated by
maximum likelihood from a historical database of international matches, with
each match contributing two Poisson log-likelihood terms; down-weighting old
matches with an exponential time decay, as advocated by Dixon and Coles
\cite{dixon1997modelling}, and shrinking $a_i, d_i$ toward zero for teams with
few observations are both recommended. The companion script
\texttt{fit\_parameters.py} implements exactly this estimation: iteratively
reweighted least squares on the 10{,}048 internationals played since 2016 in
the open dataset of \cite{martj42results} (two observations per match), with
a three-year half-life decay and a light ridge penalty on the team
parameters. Section~\ref{sec:data} reports the fitted values used throughout
this note.

The independence assumption in \eqref{eq:poisson} is a deliberate
simplification: empirically, low-scoring outcomes exhibit mild dependence,
which Dixon and Coles \cite{dixon1997modelling} capture with a local
correction factor and Karlis and Ntzoufras \cite{karlis2003analysis} with a
bivariate Poisson model. Either refinement
can be substituted for \eqref{eq:poisson} without changing anything downstream
in the simulation.

\section{Estimation of Match-Level Win Probabilities}
\label{sec:matchprob}

Under \eqref{eq:poisson}, the joint probability of the regular-time score
$(x, y)$ is the product of the two Poisson probability mass functions:
\begin{equation}
\Prob(G_i = x,\, G_j = y)
\;=\;
\frac{e^{-\lambda_{ij}}\,\lambda_{ij}^{x}}{x!}
\cdot
\frac{e^{-\lambda_{ji}}\,\lambda_{ji}^{y}}{y!},
\qquad x, y \in \{0, 1, 2, \dots\}.
\label{eq:jointpmf}
\end{equation}
Summing over the relevant regions of the score grid gives the regular-time
outcome probabilities
\begin{align}
\Prob(i \text{ wins in regular time})
&= \sum_{x > y} \Prob(G_i = x,\, G_j = y),
\label{eq:winRT}\\
\Prob(\text{draw})
&= \sum_{x = y} \Prob(G_i = x,\, G_j = y),
\label{eq:draw}
\end{align}
with $\Prob(j \text{ wins in regular time})$ the complementary sum over
$x < y$.

For numerical evaluation the infinite grid is truncated at a maximum goal
count $G_{\max}$ (we use $G_{\max} = 10$), and the truncated joint
distribution is renormalized to compensate for the discarded tail mass:
\begin{equation}
\widetilde{\Prob}(G_i = x,\, G_j = y)
\;=\;
\frac{\Prob(G_i = x,\, G_j = y)}
{\displaystyle\sum_{x'=0}^{G_{\max}}\sum_{y'=0}^{G_{\max}}
\Prob(G_i = x',\, G_j = y')},
\qquad 0 \le x, y \le G_{\max}.
\label{eq:truncation}
\end{equation}
For realistic international scoring rates ($\lambda \le 3$), the Poisson tail
mass beyond ten goals is of order $10^{-4}$ or smaller, so the truncation
error is negligible; the renormalization in \eqref{eq:truncation} removes even
that residual and guarantees that the computed probabilities sum exactly to
one. In the implementation, the sums
\eqref{eq:winRT}--\eqref{eq:draw} are evaluated on the
$(G_{\max}+1) \times (G_{\max}+1)$ outer-product grid of the two truncated
marginal probability vectors.

\section{Treatment of Draws, Extra Time, and Penalties}
\label{sec:draws}

Every remaining match is a knockout tie: if the score is level after regular
time, 30 minutes of extra time are played, followed if necessary by a penalty
shoot-out. Rather than modelling extra-time goals and penalty kicks
explicitly, we collapse the entire tie-break stage into a single conditional
probability, logistic in the Elo difference:
\begin{equation}
\Prob\bigl(i \text{ wins tie-break} \,\big|\, \text{draw}\bigr)
\;=\;
\frac{1}{1 + \exp\bigl(-\eta\,(E_i - E_j)\bigr)}.
\label{eq:tiebreak}
\end{equation}
The slope $\eta \ge 0$ controls how strongly team quality carries into extra
time and penalties: $\eta = 0$ makes the tie-break a fair coin, which is a
reasonable approximation for penalty shoot-outs, while larger $\eta$ lets the
stronger team convert its superiority. Fitted by maximum likelihood to the
681 recorded penalty shoot-outs in \cite{martj42results}
(\texttt{fit\_parameters.py}), the slope is $\eta = 0.0014$: a 100-point Elo
advantage yields a tie-break win probability of only $0.53$, confirming
empirically that shoot-outs are nearly fair coins. Because the fit uses
shoot-outs alone, it plausibly understates the stronger team's edge in the
extra-time portion of the tie-break, making this $\eta$ a conservative
choice. A more granular alternative---simulating extra time as a 30-minute
Poisson period with rates $\lambda_{ij}/3$ and modelling the shoot-out
separately---is discussed in Appendix~\ref{app:extensions};
\eqref{eq:tiebreak} is the parsimonious default.

Combining \eqref{eq:winRT}, \eqref{eq:draw}, and \eqref{eq:tiebreak}, the
\emph{total advancement probability} of team $i$ against team $j$ in a
knockout match is
\begin{equation}
p_{ij}
\;=\;
\Prob(i \text{ wins in regular time})
\;+\;
\Prob(\text{draw}) \cdot
\Prob\bigl(i \text{ wins tie-break} \,\big|\, \text{draw}\bigr).
\label{eq:advancement}
\end{equation}
Because the logistic function satisfies
$\sigma(z) + \sigma(-z) = 1$, the advancement probabilities are coherent:
$p_{ij} + p_{ji} = 1$ for every pair, so exactly one team advances.

The quantities $p_{ij}$ can be used inside the Monte Carlo simulation in two
equivalent ways. One may draw the advancing team directly as a Bernoulli
trial with success probability $p_{ij}$; or one may simulate the match at the
score level---draw $(G_i, G_j)$ from \eqref{eq:poisson} and, in case of a
draw, resolve the tie with a Bernoulli draw from \eqref{eq:tiebreak}. Both
schemes induce the same advancement law; we adopt the score-level scheme
(Algorithm~\ref{alg:match}) because it additionally yields simulated score
lines at no extra cost, while \eqref{eq:advancement} provides an analytic
cross-check on the simulator.

\section{Monte Carlo Tournament Simulation}
\label{sec:mc}

Championship probabilities compound seven dependent match outcomes: a team
must win its quarterfinal, then beat whichever team emerges from the paired
quarterfinal, then beat whichever team survives the other half of the
bracket. Monte Carlo simulation handles this compounding mechanically. In
each replication $b = 1, \dots, B$ we simulate the four quarterfinals, feed
the winners into the two semifinals, simulate the final, and record the
champion $C^{(b)}$. The estimator of \eqref{eq:target} is the empirical
frequency
\begin{equation}
\hat{p}_t \;=\; \frac{1}{B} \sum_{b=1}^{B} \Ind\{C^{(b)} = t\},
\qquad t \in \mathcal{T},
\label{eq:mcestimator}
\end{equation}
which is unbiased for $p_t$ under the model, with Monte Carlo standard error
\begin{equation}
\mathrm{se}(\hat{p}_t) \;=\; \sqrt{\frac{\hat{p}_t (1 - \hat{p}_t)}{B}}
\;\le\; \frac{1}{2\sqrt{B}}.
\label{eq:mcse}
\end{equation}
With $B = 100{,}000$ replications the standard error is at most $0.0016$
(about $0.16$ percentage points), far below any plausible model
uncertainty, and the simulation completes in about a second on a laptop. A
fixed random seed makes the results exactly reproducible.\footnote{The
simulator draws goals from the untruncated Poisson distributions in
\eqref{eq:poisson}; the truncation \eqref{eq:truncation} is used only for the
analytic probabilities \eqref{eq:advancement}. The resulting discrepancy is
below the Monte Carlo noise floor.}

With eight teams and seven matches, the champion probabilities could also be
computed \emph{exactly} by enumerating the $2^7 = 128$ bracket outcomes, or
via the closed-form path factorization given in Appendix~\ref{app:path}. We
nevertheless present the Monte Carlo route---the standard approach in
tournament forecasting from the championship simulator of Koning et
al.~\cite{koning2003simulation} to the World Cup forecasts of Groll et
al.~\cite{groll2019hybrid}: it is just as accurate for practical
purposes, its code mirrors the recursive structure of the tournament, and it
survives every extension of Appendix~\ref{app:extensions}---Bayesian parameter
draws, injury shocks, score-dependent features---for which the closed form
ceases to exist. Appendix~\ref{app:whymc} expands on this trade-off.

\section{Data Inputs}
\label{sec:data}

The model consumes one row of covariates per quarterfinalist:
Elo rating $E_i$, attack strength $a_i$, defence strength $d_i$, tournament
form $F_i$, and rest days $R_i$. All inputs are estimated from the open
match-level dataset of \cite{martj42results}, which at retrieval (9~July
2026) covered 49{,}505 internationals from 1872 through 7~July 2026,
including every match of the 2026 World Cup up to the round of~16. The
components are constructed as follows.
\begin{itemize}
  \item \textbf{Elo ratings} are computed by replaying the full match
        history through the eloratings.net update rule
        \cite{elo1978rating,eloratings}: importance-dependent $K$ factors
        (60 for World Cup finals matches, down to 20 for friendlies), a
        goal-difference multiplier, a 100-point home-advantage adjustment
        for non-neutral venues, and shoot-outs counted as draws. The
        replication reproduces the published ranking of the eight
        quarterfinalists (Spain first, then Argentina, France, England) and
        their relative gaps to within a few points, up to a benign uniform
        level offset of roughly $+70$; only rating \emph{differences} enter
        the model, so the offset is irrelevant.
  \item \textbf{Attack and defence} parameters and the coefficients
        $(\gamma, \delta, \rho)$ come from the time-decayed Poisson
        regression of Section~\ref{sec:poisson}
        (\texttt{fit\_parameters.py}), fitted jointly with a home-advantage
        term on all 2016--2026 internationals.
  \item \textbf{Form} $F_i$ is the mean goal difference over the team's last
        five internationals---for a 2026 quarterfinalist, exactly its five
        World Cup matches so far (three group games, round of~32, round
        of~16)---matching the definition used in the regression.
  \item \textbf{Rest days} $R_i$ are counted from the round-of-16 match to
        the scheduled quarterfinal (capped at 14 in the regression; all
        quarterfinalists have 4--6 days).
\end{itemize}
Table~\ref{tab:inputs} lists the resulting inputs in bracket order. The same
values are stored in \texttt{data\_template.csv}, which the simulation script
reads; re-running \texttt{fit\_parameters.py} on a refreshed download
regenerates them. Appendix~\ref{app:repro} records the full provenance
chain---download URLs, retrieval date, file checksums, every processing
rule, the fitted coefficients, and step-by-step reproduction commands.

\begin{table}[htbp]
\centering
\caption{Quarterfinal bracket and estimated model inputs (data through
7~July 2026; see the text for construction). Attack and defence are
mean-zero across all 294 teams in the regression, so top-eight values are
well above zero. Values are rounded; \texttt{data\_template.csv} carries
full precision.}
\label{tab:inputs}
\begin{tabular}{llrrrrr}
\toprule
Match & Team & Elo $E_i$ & Attack $a_i$ & Defence $d_i$ & Form $F_i$ & Rest $R_i$ (days)\\
\midrule
QF1 & France      & 2207 & 0.63 & 0.59 & 2.4 & 5\\
    & Morocco     & 2056 & 0.43 & 0.96 & 1.2 & 5\\
\addlinespace
QF2 & Spain       & 2242 & 0.73 & 0.73 & 1.8 & 4\\
    & Belgium     & 2042 & 0.69 & 0.51 & 1.6 & 4\\
\addlinespace
QF3 & Norway      & 2042 & 0.76 & 0.48 & 0.6 & 6\\
    & England     & 2150 & 0.59 & 0.77 & 1.2 & 6\\
\addlinespace
QF4 & Argentina   & 2236 & 0.58 & 0.81 & 1.8 & 4\\
    & Switzerland & 2024 & 0.61 & 0.54 & 1.2 & 4\\
\bottomrule
\end{tabular}
\end{table}

The fitted coefficients are
\[
\alpha = 0.1503,\qquad
\gamma = 0.00097,\qquad
\delta = 0.0053,\qquad
\rho = -0.0009,\qquad
\eta = 0.0014,
\]
with the score grid truncated at $G_{\max} = 10$. Three remarks aid
interpretation. First, the regression also recovers a home-advantage effect
of $\exp(0.234) \approx 1.26$---the textbook 26\% scoring boost---which
validates the pipeline but plays no role in the neutral-venue knockout
matches simulated here. Second, the fitted $\rho$ is practically nil, so the
small rest-day differences between quarterfinalists barely move the
forecasts. Third, $\gamma$ and $\delta$ are modest by construction: the
attack and defence parameters already absorb persistent quality differences,
so the Elo and form terms only capture what those parameters miss. The
resulting expected goals in the four quarterfinals range from $0.77$
(Switzerland against Argentina) to $1.76$ (Spain against Belgium), in line
with knockout-stage scoring.

\section{Algorithm}
\label{sec:algorithm}

Algorithm~\ref{alg:tournament} summarizes the Monte Carlo procedure, and
Algorithm~\ref{alg:match} the single-match subroutine that it calls seven
times per replication. The pseudocode maps one-to-one onto the functions
\texttt{simulate\_tournament} and \texttt{simulate\_match} in
\texttt{simulate\_worldcup.py}.

% ---------------------------------------------------------------------------
% algorithm.tex
% Pseudocode for the Monte Carlo tournament simulation.
% This file is \input by main.tex (Section "Algorithm"); it is not a
% standalone document.
% ---------------------------------------------------------------------------

\begin{algorithm}[htbp]
\caption{Monte Carlo estimation of conditional championship probabilities}
\label{alg:tournament}
\begin{algorithmic}[1]
\Require quarterfinal pairings $(i_k, j_k)$ for $k = 1, \dots, 4$;
         covariates $(E_t, F_t, R_t)$ and strength parameters $(a_t, d_t)$
         for the eight quarterfinalists;
         coefficients $(\alpha, \gamma, \delta, \rho, \eta)$;
         number of replications $B$
\Ensure  estimates $\hat{p}_t$ of
         $\Prob(t \text{ wins the World Cup} \mid t \text{ is a quarterfinalist})$
\State $c_t \gets 0$ for every quarterfinalist $t$
       \Comment{championship counters}
\For{$b \gets 1, \dots, B$}
    \For{$k \gets 1, \dots, 4$}
        \State $W_k \gets \Call{SimulateMatch}{i_k, j_k}$
        \Comment{quarterfinals}
    \EndFor
    \State $S_1 \gets \Call{SimulateMatch}{W_1, W_2}$
    \Comment{semifinal 1: winner QF1 vs.\ winner QF2}
    \State $S_2 \gets \Call{SimulateMatch}{W_3, W_4}$
    \Comment{semifinal 2: winner QF3 vs.\ winner QF4}
    \State $C \gets \Call{SimulateMatch}{S_1, S_2}$
    \Comment{final}
    \State $c_C \gets c_C + 1$
    \Comment{record the simulated champion}
\EndFor
\For{every quarterfinalist $t$}
    \State $\hat{p}_t \gets c_t / B$
    \Comment{title frequency estimates \eqref{eq:mcestimator}}
\EndFor
\State \Return $(\hat{p}_t)_{t \in \mathcal{T}}$
\end{algorithmic}
\end{algorithm}

\begin{algorithm}[htbp]
\caption{The \textsc{SimulateMatch} subroutine: one knockout match}
\label{alg:match}
\begin{algorithmic}[1]
\Function{SimulateMatch}{$i, j$}
    \State $\lambda_{ij} \gets \exp\bigl\{\alpha + a_i - d_j
           + \gamma (E_i - E_j) + \delta (F_i - F_j)
           + \rho (R_i - R_j)\bigr\}$
    \State $\lambda_{ji} \gets \exp\bigl\{\alpha + a_j - d_i
           + \gamma (E_j - E_i) + \delta (F_j - F_i)
           + \rho (R_j - R_i)\bigr\}$
    \State draw $G_i \sim \mathrm{Poisson}(\lambda_{ij})$ and
           $G_j \sim \mathrm{Poisson}(\lambda_{ji})$ independently
    \If{$G_i > G_j$}
        \State \Return $i$ \Comment{win in regular time}
    \ElsIf{$G_j > G_i$}
        \State \Return $j$
    \Else
        \Comment{draw after regular time: extra time and penalties}
        \State $\pi_{ij} \gets
               \bigl(1 + \exp\{-\eta\,(E_i - E_j)\}\bigr)^{-1}$
        \Comment{logistic tie-break \eqref{eq:tiebreak}}
        \State draw $U \sim \mathrm{Uniform}(0, 1)$
        \If{$U < \pi_{ij}$}
            \State \Return $i$
        \Else
            \State \Return $j$
        \EndIf
    \EndIf
\EndFunction
\end{algorithmic}
\end{algorithm}


\section{Output Interpretation}
\label{sec:output}

Running \texttt{simulate\_worldcup.py} on the estimated inputs of
Table~\ref{tab:inputs} first reports the analytic quarterfinal advancement
probabilities from \eqref{eq:advancement}---France over Morocco $0.552$,
Spain over Belgium $0.706$, England over Norway $0.608$, and Argentina over
Switzerland $0.694$---and then the simulated championship probabilities shown
in Table~\ref{tab:output}.

\begin{table}[htbp]
\centering
\caption{Conditional championship probabilities,
$\hat{p}_t = \widehat{\Prob}(\text{champion} \mid \text{quarterfinalist})$,
from $B = 100{,}000$ simulated tournaments (seed 42) under the estimated
inputs of Table~\ref{tab:inputs} (data through 7~July 2026, the eve of the
quarterfinals). Re-running the pipeline after each completed match refreshes
the forecast.}
\label{tab:output}
\begin{tabular}{lr}
\toprule
Team & $\hat{p}_t$\\
\midrule
Spain       & 0.2554\\
Argentina   & 0.2291\\
England     & 0.1451\\
France      & 0.1333\\
Morocco     & 0.0881\\
Norway      & 0.0587\\
Belgium     & 0.0480\\
Switzerland & 0.0424\\
\midrule
Total       & 1.0000\\
\bottomrule
\end{tabular}
\end{table}

Several remarks guide the reading of Table~\ref{tab:output}. The eight
probabilities sum to one because, conditional on the quarterfinal stage,
exactly one of these teams wins the title; teams eliminated earlier do not
appear and implicitly carry probability zero. Each entry carries a Monte
Carlo standard error of at most $0.0016$ by \eqref{eq:mcse}, so differences
beyond the third decimal are noise. Two features of the estimates stand out.
First, the tightest quarterfinal is France--Morocco: Morocco's fitted
defence is the strongest among the eight, which compresses France's
advancement probability to $0.55$ and makes Morocco's title probability
($0.088$) the highest of the four lower-rated teams. Second, England's
championship probability ($0.145$) exceeds France's ($0.133$) even though
France carries the higher Elo rating---a pure bracket effect, dissected in
Appendix~\ref{app:path}. The script also writes the table to
\texttt{champion\_probabilities.csv} for downstream use.

Finally, the output inherits the model's specification choices: the decay
half-life, the regression window, the ridge penalty, and the shoot-out-only
fit of $\eta$ are defensible defaults rather than uniquely correct ones, and
a sensitivity analysis over them is a sensible companion to any headline
number. The forecast is also a snapshot: when covariates or coefficients
change---a completed quarterfinal, an injury update, a re-fitted
regression---the pipeline should simply be re-run.

\section{Limitations and Extensions}
\label{sec:limitations}

The model's simplicity is deliberate, but several limitations should be kept
in mind.

\begin{itemize}
  \item \textbf{Independence of scores.} Equation~\eqref{eq:poisson} ignores
        the mild dependence between the two teams' goal counts in
        low-scoring matches; the Dixon--Coles correction
        \cite{dixon1997modelling} or the bivariate Poisson of Karlis and
        Ntzoufras \cite{karlis2003analysis} are drop-in upgrades.
  \item \textbf{Static strengths.} Attack, defence, Elo, and form are frozen
        at the quarterfinal snapshot; the model does not update them as the
        simulated tournament progresses (e.g.\ no fatigue accumulation from a
        simulated extra-time semifinal).
  \item \textbf{Collapsed tie-break.} Equation~\eqref{eq:tiebreak} folds
        extra time and penalties into one logistic step; shoot-outs in
        particular are close to coin flips, and a mis-set $\eta$ distorts the
        probabilities precisely in the tight matches that decide titles.
  \item \textbf{Estimation uncertainty.} The coefficients are maximum
        likelihood point estimates whose sampling error is not propagated
        into Table~\ref{tab:output}; the Bayesian extension of
        Appendix~\ref{app:extensions} addresses this. Two data quirks also
        matter: recorded knockout scores include extra-time goals, mildly
        inflating regular-time rates, and $\eta$ is identified from
        shoot-outs only.
  \item \textbf{Covariate quality.} Elo ratings lag sudden quality changes,
        and a form score built from at most five tournament matches is
        noisy; player-level information (injuries, suspensions) is absent.
\end{itemize}

Natural extensions---incorporating expected goals, injury adjustments,
betting-market information, and full Bayesian uncertainty
quantification---are developed in Appendix~\ref{app:extensions}. All of them
leave Algorithm~\ref{alg:tournament} untouched: they enrich either the inputs
or the parameter draws, which is exactly why the Monte Carlo architecture was
chosen.

\section{Conclusion}
\label{sec:conclusion}

We have specified a complete, reproducible pipeline for conditional World Cup
champion prediction from the quarterfinal stage: a covariate-driven Poisson
model for regular-time scores, a logistic Elo-based tie-break for knockout
draws, closed-form match advancement probabilities, and a Monte Carlo
simulation of the remaining seven matches that converts match-level
probabilities into championship probabilities for the eight surviving teams.
The accompanying pseudocode, estimation script, and simulation code make the
model runnable end to end in seconds. All inputs are fitted to the
international match record through 7~July 2026---the eve of the
quarterfinals---so Table~\ref{tab:output} is a live model forecast at the
quarterfinal stage, refreshable by re-running the pipeline as results
arrive. Beyond the 2026 World Cup, the same architecture
applies unchanged to any single-elimination bracket whose match outcomes can
be given pairwise advancement probabilities.

\appendix
% ---------------------------------------------------------------------------
% model_notes.tex
% Appendix: supplementary modelling notes.
% This file is \input by main.tex after \appendix; it is not a standalone
% document.
% ---------------------------------------------------------------------------

\section{Supplementary Model Notes}
\label{app:notes}

\subsection{Why the prediction is conditional on the quarterfinal stage}
\label{app:conditional}

A pre-tournament forecast must average over an enormous outcome space: group
compositions, group-stage results, and several knockout rounds, each of which
reshuffles who meets whom. Once the tournament reaches the quarterfinals,
almost all of that uncertainty has been \emph{resolved}. The sample space
collapses to the eight surviving teams and the seven remaining matches, and
the realized bracket---including which half each team occupies---is known
exactly. Conditioning on this information has three consequences.

First, eliminated teams are excluded by construction: a team knocked out in
the group stage or an earlier round has conditional championship probability
zero, and the model never represents it. Second, the conditional
probabilities of the eight quarterfinalists necessarily sum to one, which
makes the output directly interpretable as a complete probability
distribution over the champion. Third, the conditional forecast is generally \emph{sharper} than a
pre-tournament one: for a team that was expected to struggle, having already
reached the last eight is strong evidence of current quality, which enters the
model through the updated covariates (Elo after the earlier rounds, tournament
form, rest days). Conversely, quantities from before the tournament---title
odds, group-stage simulations---answer a different question and must not be
compared numerically to the conditional probabilities computed here.

\subsection{Why Monte Carlo simulation is useful with only seven matches
remaining}
\label{app:whymc}

With eight teams and seven binary match outcomes, the remaining tournament has
only $2^7 = 128$ possible trajectories, and the championship probabilities
admit the closed-form factorization of Appendix~\ref{app:path}. Monte Carlo
simulation is nevertheless the method of choice, for four reasons.

\begin{enumerate}
  \item \textbf{Extensibility.} The closed form exists only because match
        outcomes are conditionally independent given fixed parameters. Every
        realistic enrichment breaks it: drawing model parameters from a
        Bayesian posterior in each replication, injecting random injury or
        fatigue shocks, or letting a simulated extra-time semifinal reduce a
        finalist's rest covariate all induce dependence across matches.
        The simulator of Algorithm~\ref{alg:tournament} absorbs each of these
        changes with a few added lines; the analytic formula has to be
        re-derived or abandoned.
  \item \textbf{Richer outputs.} The same simulated trajectories that
        estimate champion probabilities also deliver, at no extra cost, the
        distribution of final pairings, the probability that a given team
        reaches each round, expected goals per fixture, and the joint law of
        any bracket event of interest.
  \item \textbf{Negligible cost and controlled error.} With $B = 100{,}000$
        replications the worst-case standard error is $0.16$ percentage
        points, far below model uncertainty, and the run completes in about a
        second. Increasing $B$ buys arbitrary accuracy at linear cost.
  \item \textbf{Transparency.} The code mirrors the tournament's structure
        line by line (four quarterfinals, two semifinals, one final), which
        makes it easy to audit against the bracket---often harder for a
        hand-derived sum over 128 paths.
\end{enumerate}

\subsection{How the bracket path affects championship probability}
\label{app:path}

Championship probability is not a function of team strength alone; it depends
on \emph{who can be met, and when}. Consider a team $t$ with quarterfinal
opponent $u$, let $\mathcal{V}$ denote the two teams in the paired
quarterfinal (the potential semifinal opponents), and let $\mathcal{W}$ denote
the four teams in the other half of the bracket (the potential final
opponents). Because match outcomes are conditionally independent given the
parameters, the conditional championship probability factorizes along the
path:
\begin{equation}
\Prob(C = t \mid \mathcal{Q})
\;=\;
\underbrace{p_{tu}}_{\text{quarterfinal}}
\;\times\;
\underbrace{\sum_{v \in \mathcal{V}} q_v\, p_{tv}}_{\text{semifinal}}
\;\times\;
\underbrace{\sum_{w \in \mathcal{W}} f_w\, p_{tw}}_{\text{final}},
\label{eq:pathfactorization}
\end{equation}
where $q_v$ is the probability that $v$ wins the paired quarterfinal, and
$f_w$ is the probability that $w$ emerges from the other half, computed
recursively as $f_w = \Prob(w \text{ wins its QF}) \sum_{u'} q_{u'}\,p_{wu'}$
with the sum running over the other half's adjacent quarterfinal (so that
$\sum_{w \in \mathcal{W}} f_w = 1$). Equation
\eqref{eq:pathfactorization} is exactly what the Monte Carlo procedure
estimates, and it makes the bracket dependence explicit: the semifinal factor
for France averages over \{Spain, Belgium\}, whereas the semifinal factor for
England averages over \{Argentina, Switzerland\}.

Under the estimated inputs this matters visibly. France enters the
quarterfinals with a higher Elo rating than England (2207 versus 2150), yet
England's championship probability (0.145) exceeds France's (0.133). The
reason is France's path: it must first break down Morocco, the best-fitted
defence among the eight, and then---because the top half's two strongest
teams, France and Spain, sit in \emph{adjacent} quarterfinals---most likely
face Spain already in semifinal~1, so at most one of the two can reach the
final. A counterfactual bracket that placed Spain in the bottom half would
raise France's title probability without changing a single team parameter.
This is why re-seeding debates are substantive: the draw itself carries
probability mass.

\subsection{Extensions: expected goals, injuries, betting odds, and Bayesian
uncertainty}
\label{app:extensions}

\paragraph{Expected goals (xG).}
Raw goal counts are noisy realizations of chance creation. A first upgrade
fits the attack and defence parameters of \eqref{eq:lograte} to xG totals
rather than goals, reducing finishing noise in estimation. A second adds an
xG-differential covariate to the log-rate,
$\log \lambda_{ij} \mathrel{+}= \kappa\,(X_i - X_j)$, where $X_i$ is the
team's tournament xG difference per match; this captures underlying
performance that results-based measures like Elo absorb only slowly.

\paragraph{Injuries and suspensions.}
Player availability can be folded in through an availability-adjusted rating
$E_i^{\mathrm{adj}} = E_i - \pi_i$, where $\pi_i$ converts missing players
into rating points (e.g.\ minutes-weighted contribution estimates). More
ambitiously, availability can be made stochastic: in each Monte Carlo
replication, draw an injury/suspension scenario for the remaining rounds and
adjust the affected team's parameters before simulating---one of the
dependence-inducing extensions for which the closed form of
Appendix~\ref{app:path} no longer applies.

\paragraph{Betting odds.}
Bookmaker prices aggregate large information sets. After removing the
bookmaker margin (overround), match odds imply advancement probabilities
$p_{ij}^{\mathrm{mkt}}$ that can be used in three ways: as a benchmark to
validate the model's $p_{ij}$; as a blend
$\tilde{p}_{ij} = \omega\, p_{ij}^{\mathrm{mkt}} + (1 - \omega)\, p_{ij}$
with $\omega$ chosen by historical calibration; or as calibration targets to
which the coefficients $(\gamma, \delta, \rho, \eta)$ are tuned.

\paragraph{Bayesian uncertainty.}
The pipeline treats $\theta = (\alpha, \{a_i\}, \{d_i\}, \gamma, \delta,
\rho, \eta)$ as known, so Table~\ref{tab:output} understates real
uncertainty. A Bayesian treatment places priors on $\theta$, obtains the
posterior $p(\theta \mid \mathcal{D})$ from historical match data
$\mathcal{D}$ (e.g.\ via MCMC in Stan or PyMC), and reports the posterior
predictive championship probability
\begin{equation}
\Prob(C = t \mid \mathcal{Q}, \mathcal{D})
\;=\;
\int \Prob(C = t \mid \mathcal{Q}, \theta)\,
p(\theta \mid \mathcal{D})\, d\theta,
\label{eq:posteriorpredictive}
\end{equation}
implemented by drawing a fresh $\theta^{(b)}$ from the posterior in each
Monte Carlo replication before simulating the bracket. The spread of
$\Prob(C = t \mid \mathcal{Q}, \theta^{(b)})$ across draws yields credible
intervals for each championship probability---typically wider, and more
honest, than the Monte Carlo error bars alone. Time-decayed likelihood
weights in the spirit of Dixon--Coles integrate naturally into the same
framework.

\section{Data Provenance and Reproducibility}
\label{app:repro}

This appendix records everything needed to audit or reproduce the numbers in
this note. All inputs derive from one public match-level dataset; one
published rating table is used solely for cross-validation.

\subsection{Sources and retrieval}
\label{app:sources}

\begin{itemize}
  \item \textbf{Match results.} \texttt{results.csv} from the
        community-maintained repository of \cite{martj42results}
        (\url{https://github.com/martj42/international_results}, raw file
        \texttt{master/results.csv}), retrieved on 9~July 2026:
        49{,}505 international matches from 30~November 1872 through
        7~July 2026, plus the four scheduled quarterfinal fixtures. The file
        contains all 96 matches of the 2026 World Cup played through the
        round of~16 (72 group matches, 16 in the round of~32, 8 in the round
        of~16). MD5 checksum:
        \texttt{3820fec01802a4303d91b00d6e298bed}.
  \item \textbf{Penalty shoot-outs.} \texttt{shootouts.csv} from the same
        repository and retrieval date: 681 shoot-outs, 1967--2026. MD5
        checksum: \texttt{3461b3d538309f7f17b63259e4b0b524}.
  \item \textbf{Cross-validation only.} The published Elo table of 7~July
        2026 \cite{eloratings} was used to verify the Elo replication of
        Section~\ref{sec:data} (ranking and pairwise gaps of the eight
        quarterfinalists); it enters no computation.
\end{itemize}
Both raw files are archived with the project sources, so every number in
this note can be reproduced offline; the checksums allow byte-level
verification against any fresh download.

\subsection{Processing rules}
\label{app:processing}

\begin{enumerate}
  \item Rows with missing scores (unplayed fixtures) are excluded from all
        estimation; matches are processed in chronological order.
  \item Recorded scores include extra-time goals and exclude penalty
        shoot-outs (dataset convention). Matches decided by shoot-out
        therefore enter both the Elo replay and the regression as draws.
  \item \emph{Elo replay:} every team starts at 1500; after each match,
        $R \leftarrow R + K G (W - W_e)$ with expected score
        $W_e = (1 + 10^{-\mathrm{dr}/400})^{-1}$, where $\mathrm{dr}$ is the
        rating difference plus 100 points for the home side at non-neutral
        venues; $K = 60$ (World Cup finals), $50$ (continental finals),
        $40$ (qualifiers, Nations League), $30$ (other tournaments), $20$
        (friendlies); goal-difference multiplier $G = 1$, $1.5$, or
        $(11 + |\Delta|)/8$ for winning margins $\le 1$, $= 2$, $\ge 3$.
  \item \emph{Poisson regression:} all matches from 1~January 2016 through
        7~July 2026 (10{,}048 matches, two observations each, 294 teams);
        time-decay weights $0.5^{\Delta t / 1095\,\mathrm{days}}$ with
        $\Delta t$ measured to 7~July 2026 (three-year half-life; effective
        sample size 7{,}891); ridge penalty of $1.0$ on the attack/defence
        dummies only; attack/defence recentred to mean zero after
        convergence; Elo differences scaled by $1/100$ inside the solver
        and rescaled on output; a home-advantage dummy is included in
        estimation (fitted $0.2342$) and switched off in the neutral-venue
        knockout predictions.
  \item \emph{Covariates:} form $F_i$ is the mean goal difference over the
        team's previous five internationals; rest $R_i$ is the number of
        days since the team's previous international, capped at 14.
        Identical definitions are used in estimation and prediction.
  \item \emph{Tie-break:} $\eta$ is fitted by Newton--Raphson maximum
        likelihood on the 681 shoot-outs, using each side's pre-match
        replicated Elo rating.
\end{enumerate}

\subsection{Reproduction steps}
\label{app:steps}

With Python $\ge 3.9$, \texttt{numpy}, and \texttt{pandas}:

\begin{verbatim}
# (optional) refresh the data; otherwise the archived copies are used
curl -sL -O https://raw.githubusercontent.com/martj42/\
international_results/master/results.csv
curl -sL -O https://raw.githubusercontent.com/martj42/\
international_results/master/shootouts.csv

python3 fit_parameters.py      # Elo replay + IRLS + eta fit  (~2 s)
python3 simulate_worldcup.py   # 100,000 simulations, seed 42 (~1 s)
\end{verbatim}

\noindent
\texttt{fit\_parameters.py} regenerates \texttt{data\_template.csv}
(Table~\ref{tab:inputs}) and prints the fitted coefficients
\[
\alpha = 0.150326,\qquad
\gamma = 0.000972,\qquad
\delta = 0.005260,\qquad
\rho = -0.000886,\qquad
\eta = 0.001382,
\]
which are hard-coded in \texttt{simulate\_worldcup.py}. Every stage of the
pipeline is deterministic---the simulation uses
\texttt{numpy.random.default\_rng(42)}---so independent runs on the archived
data reproduce Table~\ref{tab:output} exactly.


\bibliographystyle{plain}
\bibliography{references}

\end{document}
